\chapter{Metode Formal dengan Lean 4}

Sebuah berkas bernama \texttt{Sertifikat.lean} tiba bersama satu kalimat pengantar: ``semuanya
sudah dibuktikan mesin''. Isinya empat puluh teorema satu baris, hampir semuanya berakhir dengan
\texttt{by native\_decide}, dan di bagian bawah ada satu konstanta ringkasan yang menyebut ambang
$10^{-3}$. Tidak satu pun barisnya dapat kita periksa dengan mata dalam waktu yang wajar: yang
dibandingkan adalah bilangan bulat enam belas digit hasil pengubahan pecahan biner menjadi bentuk
eksak. Kalimat pengantar itu benar dalam arti yang sempit dan menyesatkan dalam arti yang penting.
Mesin memang memeriksa setiap baris, tetapi ia tidak memeriksa apakah yang diperiksa itu yang
dimaksud.

Pengujian dan pembuktian menjanjikan hal yang berbeda, dan bedanya bukan soal derajat
kepercayaan. Pengujian mengambil sejumlah masukan, menjalankan program, lalu melaporkan apa yang
terjadi pada masukan itu. Kesimpulannya selalu berbentuk ``untuk $10^{4}$ sampel uji, tidak ada
yang gagal''. Pembuktian formal berjanji lebih: satu pernyataan berkuantor atas seluruh masukan,
disertai objek yang dapat diperiksa pihak ketiga tanpa mempercayai penulisnya. Janji itu
ditepati hanya bila tiga syarat dipenuhi sekaligus, yaitu pernyataannya ditulis dengan benar,
objek yang diperiksa memang bukti pernyataan itu, dan landasan pemeriksaannya tidak disalahgunakan.
Bab ini membongkar ketiga syarat tersebut pada satu sistem konkret, Lean 4.

Alat dari bab sebelumnya belum cukup untuk pekerjaan ini. Bab 27 mengajarkan cara melaporkan
apa yang benar-benar diukur: pisahkan data latih dari data uji, sebutkan distribusinya, laporkan
variansi antar seed. Itu adalah disiplin tentang \emph{klaim empiris}. Di sini kita menghadapi
klaim yang bentuknya lain: bukan rata-rata atas sampel, melainkan pernyataan berkuantor
$\forall x,\ P(x)$. Untuk klaim semacam itu tidak ada interval kepercayaan yang menolong, karena
satu penyangkal sudah cukup untuk meruntuhkannya, dan ketiadaan penyangkal bukan bukti. Yang
dibutuhkan adalah objek yang bisa diperiksa ulang oleh orang yang tidak percaya kepada kita.

\begin{untukanak}
Bayangkan seorang juri catur yang sangat teliti. Ia tidak pernah bertanya apakah permainan itu
menarik, apakah kedua pemain bermain dengan baik, atau apakah aturan caturnya memang aturan yang
seharusnya dipakai. Tugasnya hanya satu: setiap kali ada langkah, ia memeriksa apakah langkah itu
sah menurut buku aturan. Kalau semua langkah sah, ia menandatangani papan itu dan berkata
``permainan ini sah''. Ia tidak peduli siapa yang menang.

Sekarang bayangkan seseorang menyerahkan papan yang penuh langkah sah, lalu berkata ``ini
membuktikan bahwa saya bisa mengalahkan siapa pun''. Juri tadi akan menjawab dengan jujur: yang
saya periksa hanya bahwa langkahnya sah, bukan bahwa kalimatmu benar. Kalau kalimatnya keliru,
papan yang sah tidak menolong sama sekali, karena yang salah bukan langkahnya melainkan
pertanyaannya.

Pernahkah kamu memberikan jawaban yang benar untuk soal yang salah? Atau mengerjakan pekerjaan
rumah dengan rapi padahal yang diminta halaman lain? Di bab ini kita akan tahu mengapa komputer
bisa memeriksa pembuktian sampai ke langkah terkecil, mengapa pemeriksaan itu sama sekali tidak
menjamin bahwa pertanyaannya tepat, dan bagaimana orang menulis pertanyaan yang tepat dalam
bahasa yang bisa diperiksa mesin.
\end{untukanak}

\section{Bukti informal dan bukti formal}

\tingkat{menengah}

Bukti dalam buku matematika adalah teks yang meyakinkan pembaca terlatih. Ia memuat lompatan yang
disepakati bersama: ``jelas'', ``tanpa mengurangi keumuman'', ``analog untuk kasus lain''. Lompatan
itu bukan kecurangan; ia adalah pemadatan yang diterima komunitas. Masalahnya, pemadatan semacam
itu tidak bisa diperiksa mesin, dan tidak bisa dipisahkan dari pengetahuan pembacanya.

Bukti formal adalah versi lain dari objek yang sama. Pernyataannya ditulis dalam bahasa dengan
tata bahasa dan tipe yang tegas, dan buktinya adalah sebuah \emph{term}, yaitu ekspresi yang
tipenya adalah pernyataan itu. Pemeriksaan dilakukan oleh inti sistem, sepotong kecil kode yang
hanya tahu beberapa aturan dan tidak tahu apa pun tentang matematika. Karena intinya kecil,
pemeriksaannya dapat dipercaya jauh lebih kuat daripada program besar mana pun di sekitar
sistem itu.

\kamus{verifikasi formal}{formal verification}{Pemeriksaan langkah demi langkah oleh juru tulis
yang tidak pernah bosan dan tidak pernah percaya pada penulisnya.}

\begin{definisi}[Bukti formal]
Misalkan $P$ adalah pernyataan yang ditulis dalam bahasa formal sistem. \emph{Bukti formal} $P$
adalah term $p$ sedemikian sehingga $p:P$ dapat diturunkan oleh inti sistem dari definisi,
teorema yang sudah diperiksa, dan aksioma yang dinyatakan terbuka. Sistem menerima $P$ bila dan
hanya bila term semacam itu ditemukan dan lolos pemeriksaan tipe.\index{verifikasi formal}
\end{definisi}

Definisi itu memuat tiga kata yang mudah dilewati. ``Dapat diturunkan'' berarti ada prosedur
mekanis, bukan ada orang yang setuju. ``Diperiksa'' berarti hasilnya biner: diterima atau
ditolak, tanpa derajat keyakinan. ``Dinyatakan terbuka'' berarti aksioma tercatat, bukan
tersembunyi; kita akan memakai pencatatan itu nanti sebagai alat utama.

Perbedaan yang paling sering disalahpahami adalah perbedaan antara \emph{kebenaran} dan
\emph{keterbuktian}. Lean tidak tahu apakah suatu pernyataan benar dalam dunia nyata. Ia hanya
tahu apakah ada term yang bertipe pernyataan itu. Bila pernyataannya salah, term itu tidak akan
pernah ditemukan, tetapi ketiadaan term bukan bukti bahwa pernyataan itu salah; bisa jadi kita
saja yang belum menemukannya.

\begin{bacakeras}
Suara pertama: ``Sudah saya periksa, semuanya benar.''
Suara kedua: ``Tunjukkan caranya.''
Suara pertama: ``Saya membaca ulang tiga kali.''
Suara kedua: ``Itu kebiasaan, bukan cara.''
Suara pertama: ``Saya yakin.''
Suara kedua: ``Keyakinan bukan bukti. Berikan langkahnya, satu per satu.''

Begitulah percakapan dengan mesin. Ia tidak bertanya apakah kita orang baik, apakah kita sudah
sarapan, atau apakah kita lulus dengan nilai tinggi. Ia bertanya satu hal, dengan sabar, setiap
kali, tanpa nada marah: tunjukkan langkahnya. Kalau langkahnya ada, pintu terbuka. Kalau langkahnya
tidak ada, pintu tetap tertutup, serapat sebelumnya, meskipun kita mengulang permintaan itu seratus
kali. Mesin tidak menilai niat kita; ia hanya memeriksa bentuk.

Ucapkan bersama, perlahan, lalu ulangi sekali lagi: mesin tidak percaya, ia memeriksa.
\end{bacakeras}

\section{Proposisi sebagai tipe, bukti sebagai term}

\tingkat{menengah}

Di Lean tidak ada jurang antara pernyataan dan buktinya. Keduanya term, diperiksa oleh aturan yang
sama. Setiap kalimat yang dapat bernilai benar atau salah diberi tipe $\mathtt{Prop}$; bukti
kalimat itu adalah term yang bertipe kalimat tersebut. Pernyataan ``$P$ terbukti'' karena itu
tidak berarti $P$ benar dalam suatu model, melainkan bahwa ada term bertipe $P$ yang lolos
pemeriksaan.

\kamus{proposisi}{proposition}{Kalimat yang bisa dinilai benar atau salah, seperti soal ujian yang
jawabannya harus dibuktikan, bukan sekadar dipilih.}

\begin{definisi}[Penghuni tipe]
Sebuah tipe $T$ disebut \emph{berpenghuni} bila terdapat term $t$ dengan $t:T$; notasi $t:T$
dibaca ``$t$ bertipe $T$''. Bila $P:\mathtt{Prop}$, maka $P$ disebut \emph{proposisi} dan $t:P$
disebut \emph{term bukti} bagi $P$.\index{term bukti} Dua term disebut \emph{sama secara definisi} bila keduanya
menyusut ke bentuk identik melalui perhitungan simbolik.\index{proposisi-tipe}
\end{definisi}

Kamus ini, yaitu korespondensi antara logika dan tipe, dikenal sebagai Curry--Howard.\index{Curry-Howard} Implikasi
menjadi tipe fungsi, karena bukti $P\Rightarrow Q$ harus mengubah sebarang bukti $P$ menjadi bukti
$Q$, dan itulah pekerjaan sebuah fungsi. Kuantor universal menjadi fungsi dependen, karena bukti
$\forall x,\ P(x)$ harus menghasilkan bukti $P(a)$ untuk setiap $a$ yang diberikan. Konjungsi
menjadi pasangan, karena bukti $P\wedge Q$ memuat kedua buktinya sekaligus. Baris terakhir tabel
di bawah adalah yang paling penting: $\bot$ tidak punya penghuni.

\begin{center}
\begin{tabular}{lll}
\hline
Logika & Tipe di Lean & Cara membangun term \\
\hline
$P\Rightarrow Q$   & $P\to Q$              & \texttt{fun h => ...} \\
$P\wedge Q$        & \texttt{And P Q}      & \texttt{And.intro} \\
$P\vee Q$          & \texttt{Or P Q}       & \texttt{Or.inl}, \texttt{Or.inr} \\
$\top$             & \texttt{True}         & \texttt{True.intro} \\
$\bot$             & \texttt{False}        & tidak ada penghuni \\
$\forall x,\ P(x)$ & $\forall x,\ P\,x$    & \texttt{fun x => ...} \\
$\exists x,\ P(x)$ & \texttt{Exists}       & \texttt{Exists.intro} \\
\hline
\end{tabular}
\end{center}

\kamus{term bukti}{proof term}{Struk bangunan yang bentuknya persis mengikuti cetak biru
pernyataannya; kalau bentuknya tidak cocok, bangunan ditolak.}

Term berikut adalah bukti paling kecil yang berguna: fungsi identitas sebagai bukti implikasi
identitas, dan komposisi dua fungsi sebagai bukti transitivitas implikasi. Keduanya diperiksa
tanpa memakai pengetahuan apa pun tentang isi $P$, $Q$, dan $R$.

\begin{verbatim}
-- bukti implikasi identitas, gaya term
def bukti_id {P : Prop} : P -> P := fun h => h

-- komposisi dua implikasi
def bukti_komposisi {P Q R : Prop} (h1 : P -> Q) (h2 : Q -> R) : P -> R :=
  fun hp => h2 (h1 hp)

#check bukti_id
#check bukti_komposisi
#check Prop
\end{verbatim}

Keluaran perintah terakhir, apa adanya, adalah \texttt{Prop : Type}: himpunan proposisi sendiri
adalah sebuah tipe. Dua baris sebelumnya mencetak tipe kedua definisi. Perhatikan bahwa
\texttt{\#check} tidak membuktikan apa pun; ia hanya menanyakan tipe, dan itulah yang juga
dilakukan inti sistem ketika memeriksa.

\begin{proposisi}
Untuk setiap $P:\mathtt{Prop}$ berlaku $P\to P$. Untuk $P,Q,R:\mathtt{Prop}$ berlaku
$(P\to Q)\to(Q\to R)\to P\to R$.
\end{proposisi}

\begin{proof}[Bukti]
Term $\lambda h.\,h$ bertipe $P\to P$: dalam konteks $h:P$, variabel $h$ sendiri adalah bukti $P$,
sehingga abstraksi atas $h$ menghasilkan fungsi dari $P$ ke $P$. Untuk pernyataan kedua, misalkan
$h_{1}:P\to Q$, $h_{2}:Q\to R$, dan $h_{p}:P$. Aplikasi $h_{1}\,h_{p}$ bertipe $Q$, lalu aplikasi
$h_{2}\,(h_{1}\,h_{p})$ bertipe $R$. Tiga abstraksi berturut-turut atas $h_{p}$, $h_{2}$, dan
$h_{1}$ memberi term bertipe $(P\to Q)\to(Q\to R)\to P\to R$.
\end{proof}

Pohon pada Gambar~\ref{fig:pohon-tipe} menuliskan bukti kedua itu sebagai pohon turunan: tiga daun,
dua batang aplikasi, satu akar. Daun paling kanan adalah hipotesis yang diberikan, bukan hasil
turunan.

\begin{figure}[htbp]
\centering
\begin{tikzpicture}[font=\small,
  daun/.style={draw=layergrey,fill=layergrey!6,rounded corners=2pt,inner sep=4pt},
  hip/.style={draw=layergrey,dashed,fill=layergrey!3,rounded corners=2pt,inner sep=4pt},
  batang/.style={thick,layergrey}]
  \node[daun] (h1) at (-3.7,2.7) {\texttt{h1 : P -> Q}};
  \node[daun] (hp) at (-0.6,2.7) {\texttt{hp : P}};
  \node[daun] (ap1) at (-2.2,1.35) {\texttt{h1 hp : Q}};
  \node[hip]  (h2) at (2.7,1.35) {\texttt{h2 : Q -> R}};
  \node[daun] (akar) at (0,0) {\texttt{h2 (h1 hp) : R}};
  \draw[batang] (h1.south) -- (ap1.north west);
  \draw[batang] (hp.south) -- (ap1.north east);
  \draw[batang] (ap1.south) -- (akar.north west);
  \draw[batang] (h2.south) -- (akar.north east);
\end{tikzpicture}
\caption{Pohon tipe untuk term \texttt{h2 (h1 hp)}: tiga daun dan dua batang aplikasi, dengan akar
bertipe $R$. Kotak putus-putus adalah hipotesis yang diberikan, bukan hasil turunan.}
\label{fig:pohon-tipe}
\end{figure}

Ada satu sifat yang perlu disebut karena ia menjelaskan mengapa bukti dipakai sebagai sertifikat,
bukan sebagai data. Inti Lean memperlakukan dua bukti mana pun dari proposisi yang sama sebagai
sama secara definisi. Sifat ini disebut \emph{irrelevansi bukti}, dan akibatnya sebuah bukti tidak
membawa informasi selain ``pernyataan ini benar''. Karena itu bukti dapat dibuang setelah
diperiksa, dan yang tersisa hanyalah status pernyataannya.

Satu peringatan dipasang sejak awal. Korespondensi Curry--Howard berlaku penuh untuk logika
konstruktif. Begitu kita memakai hukum tengah-tengah, $\mathtt{Classical.em}:P\vee\neg P$, kita
menambahkan aksioma, dan bukti yang memakainya berhenti bersifat komputasional: ia tetap sah,
tetapi tidak lagi dibaca sebagai program yang menghasilkan saksi. Banyak teorema pustaka besar
memuat aksioma pilihan di daftar aksiomanya, dan bagian akhir bab ini menunjukkan cara menanyakan
daftar itu langsung kepada Lean.

\section{Kuantor: aturan intro dan eliminasi}

\tingkat{menengah}

Pernyataan yang menarik hampir selalu berkuantor.\index{kuantor} ``Untuk setiap bilangan asli'' dan ``ada
bilangan asli'' bukan hiasan bahasa; keduanya mengubah jenis objek yang harus disediakan. Dalam
pandangan tipe, perbedaan itu menjadi perbedaan antara fungsi dan pasangan.

\kamus{kuantor}{quantifier}{Kata ``setiap'' dan ``ada''; yang pertama menuntut resep untuk semua
kasus, yang kedua menuntut satu contoh beserta buktinya.}

Aturan untuk kuantor universal sepenuhnya sama dengan aturan untuk fungsi. Untuk membuktikan
$\forall x,\ P(x)$, kita mengambil $x$ sebarang dan membuktikan $P(x)$ bagi $x$ itu; itulah
abstraksi. Untuk memakai hipotesis $\forall x,\ P(x)$, kita terapkan hipotesis itu pada sebuah
objek; itulah aplikasi. Aturan untuk kuantor eksistensial adalah aturan pasangan: untuk
membuktikan $\exists x,\ P(x)$ kita menyediakan satu saksi $a$ bersama bukti $P(a)$; untuk memakai
hipotesis eksistensial kita membongkar pasangan itu menjadi saksi dan buktinya.

\begin{center}
\begin{tabular}{lll}
\hline
Bentuk pernyataan & Cara membuktikan & Cara memakai \\
\hline
$\forall x,\ P(x)$      & \texttt{intro x}              & terapkan pada objek \\
$\exists x,\ P(x)$      & \texttt{exact Exists.intro a h} & \texttt{cases} pada hipotesis \\
$P\wedge Q$             & \texttt{constructor}          & \texttt{h.1}, \texttt{h.2} \\
$P\vee Q$               & \texttt{left}, \texttt{right} & \texttt{cases} pada hipotesis \\
$\neg P$                & \texttt{intro h} lalu capai $\bot$ & terapkan pada bukti $P$ \\
\hline
\end{tabular}
\end{center}

\begin{tangan}
Siapkan sepuluh kartu kosong dan satu spidol. Di lima kartu tulis aturan masuk: ``ambil sebarang'',
``sediakan saksi'', ``pasangkan dua bukti'', ``pilih kiri'', ``pilih kanan''. Di lima kartu lain
tulis aturan pakai: ``terapkan pada objek'', ``bongkar pasangan'', ``ambil komponen pertama'',
``pecah per kasus'', ``pakai pada bukti lawan''.

Letakkan satu kartu tujuan di ujung meja, misalnya ``setiap bilangan asli nol atau penerus''. Susun
rantai ke arah kartu tujuan: kartu pertama harus aturan masuk yang bentuknya cocok dengan tujuan,
bukan aturan pakai. Setiap kartu berikutnya hanya boleh bekerja pada sisa yang masih terbuka. Kartu
yang tidak punya sisa untuk dikerjakan disingkirkan ke tepi meja; di Lean kartu itu tidak akan
pernah dipakai.

Yang harus kamu lihat: rantai memendek menuju tujuan, dan tidak pernah ada dua kartu yang mengerjakan
sisa yang sama. Kalau itu terjadi, rantaimu bercabang, dan mesin menuntut dua bukti, bukan satu.
\end{tangan}

Contoh berikut menunjukkan keduanya bekerja pada bilangan asli. Pernyataan pertama adalah
tautologi yang tidak memerlukan penalaran apa pun; pernyataan kedua menyatakan bahwa setiap
bilangan asli adalah nol atau penerus sesuatu, dan pembuktiannya memerlukan pemecahan kasus pada
bentuk bilangan.

\begin{verbatim}
-- kuantor universal: ambil n sebarang, lalu hitung
theorem nol_kanan : forall n : Nat, n + 0 = n := by
  intro n
  rfl

-- kuantor eksistensial: sediakan saksi, lalu buktikan sifatnya
theorem ada_tiga : exists n : Nat, n = 3 := by
  exact Exists.intro 3 rfl

-- setiap bilangan asli adalah nol atau penerus sesuatu
theorem nol_atau_suksesor (n : Nat) : n = 0 \/ exists m : Nat, n = m + 1 := by
  cases n with
  | zero => left; rfl
  | succ k => right; exact Exists.intro k rfl
\end{verbatim}

Ketiga teorema itu diperiksa tanpa aksioma tambahan, dan itu dapat dipastikan dengan perintah
\texttt{\#print axioms}, yang akan kita pakai berulang kali. Perhatikan juga bahwa
\texttt{Exists.intro k rfl} memakai \texttt{rfl} sebagai bukti kesamaan: pada kasus
\texttt{succ k}, tujuannya adalah $k+1=k+1$, yang benar secara definisi.

\begin{catatan}
Arah sebaliknya tidak berlaku otomatis. Dari $\exists x,\ P(x)$ kita tidak boleh menyimpulkan
$\forall x,\ P(x)$, dan Lean akan menolak setiap percobaan ke arah itu. Kesalahan menukar urutan
kuantor adalah kesalahan spesifikasi yang paling sering terjadi, dan ia lolos tanpa keluhan dari
pengujian karena semua sampel uji memenuhinya. Bagian tentang batas di akhir bab ini membahas
akibatnya.
\end{catatan}

\section{Taktik: mengubah tujuan menjadi term}

\tingkat{dasar}

Menulis term bukti secara langsung cepat menjadi tidak terbaca, karena term yang dalam harus
dibaca dari dalam ke luar. Karena itu Lean menyediakan mode taktik.\index{taktik} Setelah kata \texttt{by},
kita tidak menulis term, melainkan urutan perintah yang mengubah \emph{keadaan tujuan}. Keadaan
tujuan dicetak sebagai daftar hipotesis, lalu garis pemisah, lalu target.

\kamus{taktik}{tactic}{Perintah memasak langkah demi langkah, bukan resep lengkap yang harus
dihafalkan sekaligus.}

\begin{verbatim}
-- keadaan tujuan dicetak setelah setiap taktik; baris "|-" dibaca "maka"
theorem imp_trans {P Q R : Prop} (h1 : P -> Q) (h2 : Q -> R) : P -> R := by
  intro hp
  trace_state          -- P Q R : Prop, h1 : P -> Q, h2 : Q -> R, hp : P, |- R
  apply h2
  trace_state          -- target berubah menjadi |- Q
  exact h1 hp
\end{verbatim}

Keluaran \texttt{trace\_state} pada mesin uji ini adalah dua blok berikut. Lean mencetak panah
sebagai karakter Unicode; di sini semuanya ditulis ulang dalam ASCII agar konsisten dengan seluruh
bab.

\begin{verbatim}
P Q R : Prop
h1 : P -> Q
h2 : Q -> R
hp : P
|- R

P Q R : Prop
h1 : P -> Q
h2 : Q -> R
hp : P
|- Q
\end{verbatim}

Taktik \texttt{intro} mengambil target berbentuk implikasi dan memindahkan antesedennya menjadi
hipotesis; itu persis abstraksi fungsi. Taktik \texttt{exact} menutup target dengan term yang
tipenya cocok persis. Taktik \texttt{apply} menyamakan kesimpulan sebuah lemma dengan target, lalu
meninggalkan sisa argumen sebagai tujuan baru. Tiga taktik ini sudah cukup untuk semua bukti
logika murni di bab ini, dan ketiganya adalah bentuk permukaan dari aturan yang sama.

Gambar~\ref{fig:dua-keadaan} memadatkan dua keadaan itu menjadi dua kotak. Yang berubah hanya baris
di bawah garis: sasaran menyusut dari $R$ ke $Q$. Daftar hipotesis tidak bertambah satu baris pun,
karena \texttt{apply} tidak menciptakan bukti baru; ia hanya menurunkan tuntutan.

\begin{figure}[htbp]
\centering
\begin{tikzpicture}[font=\small,
  kotak/.style={draw=layergrey,fill=layergrey!6,rounded corners=3pt,align=left,inner sep=6pt,
                text width=3.9cm},
  panah/.style={-{Stealth[length=2.2mm]},thick,layerteal}]
  \node[kotak] (awal) at (0,0) {\textbf{keadaan 1}\\[3pt]
    \texttt{P Q R : Prop}\\
    \texttt{h1 : P -> Q}\\
    \texttt{h2 : Q -> R}\\
    \texttt{hp : P}\\[3pt]
    \rule{3.6cm}{0.4pt}\\
    $\vdash$ \texttt{R}};
  \node[kotak] (akhir) at (6.7,0) {\textbf{keadaan 2}\\[3pt]
    \texttt{P Q R : Prop}\\
    \texttt{h1 : P -> Q}\\
    \texttt{h2 : Q -> R}\\
    \texttt{hp : P}\\[3pt]
    \rule{3.6cm}{0.4pt}\\
    $\vdash$ \texttt{Q}};
  \draw[panah] (awal.east) -- node[above,font=\scriptsize]{\texttt{apply h2}} (akhir.west);
\end{tikzpicture}
\caption{Dua keadaan tujuan pada bukti \texttt{imp\_trans}: hipotesis tetap empat baris, sasaran
menyusut dari $R$ menjadi $Q$ setelah satu taktik. Kotak kedua masih terbuka; ia ditutup
\texttt{exact h1 hp}.}
\label{fig:dua-keadaan}
\end{figure}

\begin{center}
\begin{tabular}{lll}
\hline
Taktik & Yang dikerjakan & Kapan dipakai \\
\hline
\texttt{intro}     & anteseden menjadi hipotesis   & target berbentuk $P\to Q$ \\
\texttt{exact}     & menutup target dengan term    & term sudah tersedia \\
\texttt{apply}     & kesimpulan lemma cocok target & lemma berujung sama \\
\texttt{rw}        & menulis ulang dengan persamaan & ada hipotesis $a=b$ \\
\texttt{simp}      & menyederhanakan berulang      & banyak langkah sepele \\
\texttt{cases}     & memecah target per konstruktor & perlu analisis kasus \\
\texttt{induction} & basis dan langkah induksi     & target rekursif atas \texttt{Nat} \\
\hline
\end{tabular}
\end{center}

Taktik \texttt{rw} menulis ulang target memakai persamaan. Ia bekerja pada arah yang diberikan,
dan arah itu penting: menulis ulang ke arah yang salah tidak akan menemukan pola yang dicari.
Perhatikan dua percobaan berikut, yang satu berhasil dan yang satu ditolak.

\begin{verbatim}
theorem geser (a b : Nat) (h : a = b) : a + 1 = b + 1 := by
  rw [h]        -- target menjadi b + 1 = b + 1, lalu ditutup secara definisi
  rfl

-- komutatifitas penjumlahan TIDAK berlaku secara definisi, jadi rfl gagal
theorem tukar_salah (n : Nat) : n + 1 = 1 + n := by
  rfl           -- error: The rfl tactic failed

-- versi yang benar: menulis ulang dengan lemma komutatifitas
theorem tukar_benar (n : Nat) : n + 1 = 1 + n := by
  rw [Nat.add_comm]
\end{verbatim}

Kegagalan di tengah itu adalah pelajaran, bukan kecelakaan. Penjumlahan pada bilangan asli
didefinisikan dengan rekursi pada argumen kedua, sehingga $n+1$ menyusut menjadi
$\mathtt{succ}\,n$ sedangkan $1+n$ menyusut menjadi $\mathtt{succ}\,(0+n)$, dan keduanya berhenti
di bentuk yang berbeda. Yang kita sebut ``sifat komutatif'' di sekolah adalah teorema, bukan
definisi, dan Lean menolak memperlakukannya seolah-olah definisi. Inilah beda paling terasa antara
membaca dan menulis bukti formal: setiap hal yang biasa kita lewati harus disebutkan namanya.

Taktik \texttt{cases} memecah target menurut konstruktor tipe. Pada bilangan asli ia memberi dua
kasus, nol dan penerus. Taktik \texttt{simp} menutup tujuan dengan menulis ulang memakai kumpulan
lemma yang ditandai; ia berguna tetapi juga berbahaya, karena keluarannya bergantung pada
kumpulan lemma yang sedang aktif. Bukti yang diawali \texttt{simp} tanpa argumen bisa berhenti
bekerja ketika pustaka berubah, dan itulah sebabnya pada bab ini \texttt{simp} selalu disebut
bersama daftar lemma yang dipakai.

Satu peringatan tentang taktik yang tidak ada di sini. Nama-nama seperti \texttt{linarith},
\texttt{ring}, \texttt{norm\_num}, dan \texttt{induction'} berasal dari pustaka Mathlib, bukan
dari inti Lean. Semua potongan kode di bab ini dijalankan dengan Lean 4.10.0 tanpa Mathlib,
sehingga taktik itu tidak tersedia dan tidak dipakai; sebagai gantinya dipakai \texttt{omega}
(yang ada di inti) dan \texttt{simp} dengan daftar lemma eksplisit. Perbedaan ini bukan soal
gaya: pustaka menambah ribuan lemma dan beberapa taktik, dan setiap tambahan itu ikut menentukan
apa yang bisa ditutup dalam satu baris.

\begin{untukanak}
Ambil setumpuk batu kecil. Satu batu kamu susun menjadi bujur sangkar terkecil. Lalu tambahkan
tiga batu mengelilinginya, dan hasilnya bujur sangkar yang lebih besar. Tambahkan lima batu lagi,
dan hasilnya bujur sangkar yang lebih besar lagi. Tambahkan tujuh, sembilan, sebelas, dan
seterusnya, selalu bertambah dua, dan setiap kali hasilnya tetap bujur sangkar yang rapi.

Sekarang pertanyaannya: apakah ini akan terus begitu kalau kita menambahkan sampai seratus kali,
sampai sejuta kali? Kamu tidak bisa mencoba semuanya. Kamu juga tidak perlu. Kalau kamu bisa
menunjukkan bahwa pola penambahan batu selalu menutup satu sudut dan satu sisi bujur sangkar
dengan rapi, maka satu penjelasan kecil berlaku untuk semua ukuran, termasuk ukuran yang tidak
akan pernah bisa kamu susun dengan tangan.

Komputer pemalas tidak mau mencoba sejuta kasus, dan komputer yang terlalu percaya diri tidak mau
memeriksa apa pun. Di bab ini kita akan tahu mengapa ada cara ketiga: menjelaskan polanya sekali,
lalu meminta mesin memeriksa penjelasan itu sampai ke huruf terakhirnya.
\end{untukanak}

\section{Bukti induksi lengkap: jumlah bilangan ganjil}

\tingkat{lanjut}

Kita akan membuktikan satu pernyataan aritmetika yang seluruhnya dapat diperiksa mesin, yaitu
jumlah $n$ bilangan ganjil pertama sama dengan $n^{2}$. Pernyataan ini dipilih karena tiga alasan:
rumusnya dapat dihitung dengan Python tanpa Lean, buktinya memerlukan induksi, dan langkah
aljabarnya cukup kecil sehingga setiap taktik dapat dipertanggungjawabkan.

\kamus{induksi}{induction}{Menjelaskan pola dengan dua hal saja: kasus terkecil, dan cara naik
satu langkah dari kasus mana pun.}

Sebelum menulis kode, kita nyatakan pernyataannya dan buktinya dalam bahasa biasa. Fungsinya
didefinisikan secara rekursif: jumlah nol bilangan ganjil pertama adalah nol, dan jumlah $n+1$
bilangan ganjil pertama adalah jumlah $n$ bilangan ganjil pertama ditambah bilangan ganjil ke
$n+1$, yaitu $2n+1$.

\begin{definisi}[Jumlah bilangan ganjil]
Fungsi $S:\mathbb{N}\to\mathbb{N}$ didefinisikan oleh $S(0)=0$ dan
$S(n+1)=S(n)+(2n+1)$ untuk setiap $n\in\mathbb{N}$.\index{induksi}
\end{definisi}

Angka $2n+1$ bukan pilihan sembarang: bilangan ganjil ke-$n$ memang berbentuk itu, dan
$2n+1$ juga selisih antara $(n+1)^{2}$ dan $n^{2}$. Kesamaan terakhir itulah inti buktinya.

\begin{lemma}
Untuk setiap $k\in\mathbb{N}$ berlaku $(k+1)^{2}=k^{2}+2k+1$.
\end{lemma}

\begin{proof}[Bukti]
$(k+1)^{2}=(k+1)(k+1)=k(k+1)+(k+1)=k^{2}+k+k+1=k^{2}+2k+1$. Setiap langkah memakai
distributivitas dan asosiativitas penjumlahan pada bilangan asli.
\end{proof}

\begin{teorema}
Untuk setiap $n\in\mathbb{N}$ berlaku $S(n)=n^{2}$.
\end{teorema}

\begin{proof}[Bukti]
Dengan induksi pada $n$. Untuk $n=0$, kedua sisi bernilai $0$. Misalkan pernyataan benar untuk
$k$, yaitu $S(k)=k^{2}$. Maka
$S(k+1)=S(k)+(2k+1)=k^{2}+2k+1=(k+1)^{2}$, dengan lemma di atas pada langkah terakhir. Jadi
pernyataan benar untuk $k+1$. Menurut prinsip induksi, pernyataan benar untuk setiap
$n\in\mathbb{N}$.
\end{proof}

Bukti di atas adalah bukti informal yang lengkap, dalam arti tidak ada langkah yang dilewati.
Sekarang kita menuliskannya untuk mesin. Perhatikan bahwa definisi rekursif dan teoremanya hampir
sama kata demi kata dengan yang baru kita tulis, sedangkan buktinya berubah menjadi urutan taktik.

\begin{verbatim}
-- definisi rekursif: dua persamaan, tanpa ":=" pada baris pertama
def oddSum : Nat -> Nat
  | 0     => 0
  | n + 1 => oddSum n + (2 * n + 1)

theorem oddSum_sq (n : Nat) : oddSum n = n * n := by
  induction n with
  | zero => rfl
  | succ k ih =>
    rw [oddSum]        -- buka definisi: oddSum (k+1) = oddSum k + (2*k+1)
    rw [ih]            -- pakai hipotesis induksi: oddSum k = k*k
    simp [Nat.succ_eq_add_one, Nat.mul_succ, Nat.add_mul, Nat.one_mul]
    omega              -- tutup sisa linear
\end{verbatim}

Baris \texttt{induction n with} membangkitkan dua tujuan. Tujuan pertama bertanda \texttt{zero}
dan berbunyi $S(0)=0\cdot 0$, yang ditutup \texttt{rfl} karena kedua sisi menyusut ke nol. Tujuan
kedua bertanda \texttt{succ k} dan membawa satu hipotesis tambahan bernama \texttt{ih}, yaitu
$S(k)=k\cdot k$. Nama \texttt{ih} dipilih sendiri; yang penting perannya.

Baris \texttt{rw [oddSum]} membuka definisi pada sisi kiri. Setelah baris itu, tujuan berbunyi
$S(k)+(2k+1)=(k+1)(k+1)$. Baris \texttt{rw [ih]} mengganti $S(k)$ dengan $k\cdot k$, sehingga
tujuan menjadi $k\cdot k+(2k+1)=(k+1)(k+1)$.

Di titik ini kita berada persis di tempat lemma tadi bekerja, tetapi mesin belum tahu bahwa
$(k+1)(k+1)$ dapat ditulis sebagai $k\cdot k+2k+1$. \texttt{simp} dengan empat lemma eksplisit
mengubah bentuk ruas kanan menjadi $k\cdot k+k+(k+1)$, yaitu bentuk yang seluruhnya linear dalam
$k$. Baris terakhir, \texttt{omega}, menutup tujuan linear semacam itu.

Perlu ditegaskan bahwa \texttt{omega} bukan sulap dan bukan pula aksioma tersembunyi. Ia adalah
prosedur keputusan untuk aritmetika linear atas bilangan bulat, dan ia hanya menerima tujuan yang
seluruh operasinya linear. Bila kita memanggilnya sebelum \texttt{simp}, tujuannya masih memuat
$k\cdot k$ dan $(k+1)(k+1)$ sekaligus. Pada mesin uji ini, berkas percobaan sementara bernama
\texttt{odd.lean} keluarannya menolak:

\begin{verbatim}
odd.lean:11:4: error: omega could not prove the goal:
a possible counterexample may satisfy the constraints
  a - b + 2*c <= -2
  ...
\end{verbatim}

Pesan itu jujur dan berguna: ia menyebut kemungkinan penyangkal, dan justru karena itu kita tahu
bahwa yang gagal adalah bentuk tujuan, bukan kebenaran pernyataannya. Urutan taktik karena itu
bukan selera; ia mengubah jenis tujuan yang dihadapi prosedur keputusan.

\begin{catatan}
Bila \texttt{simp} di atas dijalankan tanpa \texttt{omega}, tujuan yang tersisa berbentuk
$k\cdot k+(1+2k)=k+(k+(k\cdot k+1))$. Secara matematis itu benar, tetapi tidak linear lagi, dan
penyederhana tidak menutupnya. Pada mesin uji ini percobaan tersebut berakhir dengan
\texttt{unsolved goals}. Inilah bentuk paling umum dari kegagalan bukti formal: bukan
pernyataannya salah, melainkan bentuk tujuannya tidak cocok dengan alat yang dipakai.
\end{catatan}

\begin{akibat}
Untuk setiap $n\in\mathbb{N}$, jumlah $n$ bilangan ganjil pertama adalah kuadrat sempurna.
\end{akibat}

Akibat itu hanya pengulangan teorema dengan kata lain, dan itu memang gunanya: ia menyatakan
kembali hasil dalam bahasa yang dipakai di bagian ini, sehingga pembaca yang menemui rumusnya di
tempat lain dapat mengenalinya.

\section{Struktur, typeclass, dan apa yang benar-benar dijamin}

\tingkat{lanjut}

Sertifikat jarang berupa satu teorema lepas. Biasanya ia sekumpulan data bersama klaim tentang
data itu: ambang, tabel, nama berkas, dan beberapa pernyataan. Lean menyediakan dua alat untuk
mengelompokkan hal semacam ini, yaitu \emph{structure} dan \emph{typeclass}.

\kamus{typeclass}{type class}{Daftar kemampuan yang harus dimiliki sebuah tipe agar bisa dipakai
di tempat tertentu, seperti syarat lisensi untuk mengemudikan kendaraan jenis tertentu.}

\begin{definisi}[Struktur dan typeclass]
Sebuah \emph{struktur} adalah tipe yang bidangnya diberi nama, sehingga nilainya dibangun dengan
menyebut nama bidang. Sebuah \emph{typeclass} adalah struktur yang dipakai untuk mencari
\emph{instance} secara otomatis: ketika sebuah definisi memerlukan kemampuan tertentu dari sebuah
tipe, Lean mencari instance yang cocok di lingkungan yang sedang aktif.\index{typeclass}
\end{definisi}

Perbedaan praktisnya sederhana. Struktur dipakai ketika data itu sendiri yang penting, misalnya
sertifikat dengan bidang klaim dan ambang. Typeclass dipakai ketika yang penting adalah kemampuan:
sesuatu yang dapat dipakai sebagai bukti, sesuatu yang dapat dihitung normanya, sesuatu yang dapat
dibandingkan. Pada keduanya, pemeriksaan tetap sama ketatnya: setiap bidang harus diisi, dan
setiap instance harus lolos pemeriksaan tipe.

\begin{verbatim}
-- struktur: data yang bidangnya bernama
structure Sertifikat where
  klaim  : String
  ambang : Nat

-- typeclass: kemampuan yang dicari otomatis
class Terperiksa (a : Type) where
  jumlah_teorema : a -> Nat

instance : Terperiksa Sertifikat where
  jumlah_teorema _ := 1

def contoh : Sertifikat := { klaim := "oddSum", ambang := 100 }

#eval (Terperiksa.jumlah_teorema contoh)   -- mencetak 1
\end{verbatim}

Sekarang bagian yang paling sering dilewati dalam pembicaraan populer tentang Lean, yaitu
pertanyaan apa yang sebenarnya dijamin oleh penerimaan sebuah teorema. Empat hal dijamin, dan
empat hal lain tidak.

Gambar~\ref{fig:alur-periksa} menaruh ketiga tahap itu pada satu baris: pernyataan ditulis, term
diajukan, inti memutuskan. Hanya ada dua keluaran; tidak ada kotak ketiga yang berbunyi ``mungkin''.

\begin{figure}[htbp]
\centering
\begin{tikzpicture}[font=\small,
  tahap/.style={draw=layergrey,fill=layergrey!6,rounded corners=3pt,align=center,inner sep=5pt,
                text width=2.9cm,minimum height=1.7cm},
  panah/.style={-{Stealth[length=2.2mm]},thick,layergrey},
  terima/.style={draw=layerjade,fill=layerjade!8,rounded corners=3pt,align=center,inner sep=5pt,
                 text width=2.9cm},
  tolak/.style={draw=accent,fill=accent!6,rounded corners=3pt,align=center,inner sep=5pt,
                text width=2.9cm}]
  \node[tahap] (s1) at (0,0) {\textbf{1. pernyataan}\\$P$ bertipe \texttt{Prop}};
  \node[tahap] (s2) at (3.9,0) {\textbf{2. term}\\$p$ diajukan untuk $P$};
  \node[tahap] (s3) at (7.8,0) {\textbf{3. inti}\\periksa apakah $p:P$};
  \draw[panah] (s1) -- (s2);
  \draw[panah] (s2) -- (s3);
  \node[terima] (ya) at (5.4,-2.9) {\textbf{diterima}\\ada term bertipe $P$};
  \node[tolak] (tidak) at (9.9,-2.9) {\textbf{ditolak}\\tidak ada term};
  \draw[panah] (s3.south west) -- node[above left,font=\scriptsize]{cocok} (ya.north);
  \draw[panah] (s3.south east) -- node[above right,font=\scriptsize]{tidak cocok} (tidak.north);
\end{tikzpicture}
\caption{Alur pemeriksaan: tiga tahap (pernyataan, term, inti) dan dua keluaran, diterima atau
ditolak. Tidak ada keluaran ketiga; berkas yang berhenti di tengah dilaporkan sebagai tujuan yang
belum tertutup, bukan sebagai ``hampir benar''.}
\label{fig:alur-periksa}
\end{figure}

Yang dijamin: (i) ada term yang tipenya persis pernyataan yang ditulis; (ii) term itu diperiksa
inti, bukan oleh taktik yang menghasilkannya; (iii) pemeriksaan hanya memakai definisi, teorema
lain yang sudah diperiksa, dan aksioma yang tercatat; (iv) daftar aksioma yang dipakai dapat
ditanyakan kapan saja dengan \texttt{\#print axioms}. Butir (iv) adalah alat paling berguna di
seluruh bab ini, karena ia mengubah ``percaya saja'' menjadi ``periksa sendiri''.

Yang tidak dijamin: (i) bahwa pernyataannya menyatakan apa yang kita maksud; (ii) bahwa
definisinya memodelkan benda di dunia nyata; (iii) bahwa inti, pengurai, dan perangkat kerasnya
bebas dari kesalahan; (iv) bahwa taktik yang dipakai tidak menambahkan asumsi tersembunyi.
Butir (iv) bukan cela teoretis: ia dapat diperiksa, dan hasilnya sering mengejutkan.

\begin{contoh}[Dua cara membuktikan hal yang sama, dua daftar aksioma]
Perhatikan dua pembuktian pernyataan $S(10)=100$ yang keduanya benar. Yang pertama memakai
\texttt{decide}, yang menjalankan perhitungan di dalam inti. Yang kedua memakai
\texttt{native\_decide}\index{native-decide}, yang mengompilasi perhitungan menjadi kode mesin
dan menjalankannya. Kedua pernyataannya sama, tetapi daftar aksiomanya berbeda, dan perbedaan itu
terlihat dari perintah \texttt{\#print axioms}.
\end{contoh}

\begin{verbatim}
def oddSum : Nat -> Nat
  | 0     => 0
  | n + 1 => oddSum n + (2 * n + 1)

theorem cek_native : oddSum 10 = 100 := by native_decide
theorem cek_kernel : oddSum 10 = 100 := by decide

#print axioms cek_native
#print axioms cek_kernel
\end{verbatim}

Keluaran pada mesin uji ini, apa adanya:

\begin{verbatim}
'cek_native' depends on axioms: [Lean.ofReduceBool]
'cek_kernel' does not depend on any axioms
\end{verbatim}

Baris pertama adalah inti dari seluruh pembahasan tentang \texttt{native\_decide}. Perhitungannya
tidak dilakukan oleh inti, melainkan oleh kode terkompilasi, dan Lean mencatat kepercayaan itu
sebagai aksioma bernama \texttt{Lean.ofReduceBool}. Secara praktis ini bukan bencana: kompilator
Lean dipakai untuk mengompilasi Lean sendiri setiap hari. Secara prinsipil ini perbedaan besar,
karena batas kepercayaan berpindah dari inti yang kecil ke rantai panjang kompilator, penghubung,
dan prosesor. Sertifikat yang dibangun di atas \texttt{native\_decide} karena itu lebih lemah
daripada sertifikat yang dibangun di atas \texttt{decide}, meskipun keduanya mencetak ``diterima''
tanpa keluhan.

Pada teorema induksi tadi, daftar aksiomanya juga tidak kosong. Perintah \texttt{\#print axioms}
untuk \texttt{oddSum\_sq} mencetak \texttt{[propext, Quot.sound]}. Keduanya adalah aksioma standar yang dipakai untuk
penalaran klasik dan hasil bagi, dan keduanya muncul karena \texttt{omega} membuktikan tujuan
dengan kontradiksi di dalam bilangan bulat. Ini bukan kesalahan; ini fakta yang seharusnya
disebutkan bila sertifikat dipakai sebagai dasar klaim.

\section{Batas: menolak bukti salah, \texttt{sorry}, dan aksioma}

\tingkat{lanjut}

Kemampuan menolak sama pentingnya dengan kemampuan menerima. Berkas berikut sengaja berisi
kesalahan jenis yang berbeda, dan keluarannya adalah jawaban Lean untuk masing-masing.

\begin{verbatim}
-- Berkas: Salah.lean
theorem salah_arit : 2 + 2 = 5 := by decide

theorem salah_rfl (n : Nat) : n + 1 = 1 + n := by rfl

axiom magic : 1 = 2
theorem from_magic : 1 = 2 := magic
#print axioms from_magic

theorem belum_selesai (n : Nat) : n + 0 = n := by sorry
\end{verbatim}

Keluaran sebenarnya, dipotong pada bagian yang panjang:

\begin{verbatim}
Salah.lean:2:37: error: tactic 'decide' proved that the proposition
  2 + 2 = 5
is false
Salah.lean:4:50: error: The rfl tactic failed. Possible reasons:
- The goal is not a reflexive relation ...
- The arguments of the relation are not equal.
...
n : Nat
|- n + 1 = 1 + n
'from_magic' depends on axioms: [magic]
Salah.lean:10:8: warning: declaration uses 'sorry'
\end{verbatim}

Tiga jenis penolakan muncul di sini, dan ketiganya berbeda sifat. Yang pertama adalah penolakan
paling kuat: \texttt{decide} bukan hanya gagal menutup tujuan, ia membuktikan bahwa proposisinya
salah. Untuk pernyataan berhingga, kegagalan dan penyangkalan adalah hal yang sama. Yang kedua
adalah penolakan yang lebih halus: \texttt{rfl} menolak karena kedua sisi tidak sama secara
definisi, padahal pernyataannya benar. Di sini Lean tidak menyalahkan pernyataannya, ia
menyalahkan alat yang dipakai; pesannya bahkan menyarankan lemma refleksivitas eksplisit. Yang
ketiga bukan penolakan sama sekali.

Baris \texttt{axiom magic : 1 = 2} diterima tanpa keluhan. Lean tidak memeriksa kewajaran isi
aksioma; ia hanya mencatatnya.\index{aksioma} Setelah itu, \texttt{from\_magic} juga diterima, dan seluruh
aritmetika bisa dibongkar dengan \texttt{False.elim} dari sana. Satu-satunya tanda bahaya adalah
baris \texttt{\#print axioms}, yang menyebut \texttt{[magic]} dengan tenang, tanpa peringatan
dan tanpa warna merah. Inilah alasan praktis mengapa memeriksa daftar aksioma adalah kebiasaan
yang wajib: sistem tidak akan menegur kita.

\kamus{aksioma}{axiom}{Aturan yang kita minta diterima tanpa bukti, seperti asumsi pembuka dalam
perdebatan; sah, tetapi harus disebut.}

Baris terakhir berkas itu memakai \texttt{sorry}, dan keluarannya hanya peringatan
\texttt{declaration uses 'sorry'}.\index{sorry} Berkas tetap diproses, teorema tetap ``ada'', dan siapa pun
yang membaca keluaran dengan saksama akan melihat kata \texttt{warning}, bukan \texttt{error}.
Pada pustaka besar, satu \texttt{sorry} yang tertinggal di tengah ribuan teorema nyaris tidak
terlihat. Lean menyediakan jawabannya, dan jawabannya sama dengan jawaban untuk aksioma:
\texttt{\#print axioms} mencetak \texttt{[sorryAx]}. Pustaka matematika Lean, yaitu Mathlib,
memakai kebijakan yang lebih tegas: berkas yang memuat \texttt{sorry} ditolak oleh pemeriksaan
otomatis sebelum diterima. Kebijakan itu pantas ditiru siapa pun yang menerbitkan sertifikat.

\kamus{sorry}{sorry}{Catatan ``nanti saya kerjakan'' yang membuat pekerjaan tampak selesai;
mesin menerimanya, tetapi mencatatnya sebagai utang.}

Sekarang batas yang paling penting, dan yang paling jarang disebut. Verifikasi formal menjamin
bahwa term cocok dengan pernyataan. Ia tidak menjamin bahwa pernyataan cocok dengan maksud. Bila
spesifikasinya terlalu lemah, buktinya bisa sempurna dan tidak berguna.\index{spesifikasi} Contoh berikut
menunjukkannya dalam bentuk paling telanjang: pernyataan pertama mengatakan ``ada gangguan kecil
yang tidak mengubah keluaran'', pernyataan kedua mengatakan ``setiap gangguan kecil tidak
mengubah keluaran''. Yang pertama dapat dibuktikan untuk \emph{semua} fungsi, termasuk fungsi
terburuk yang bisa dibayangkan, karena saksi $\delta=0$ selalu tersedia.

\begin{verbatim}
def identitas (n : Nat) : Nat := n

-- spesifikasi lemah: ADA delta yang tidak mengubah keluaran (untuk semua f)
theorem ada_delta (f : Nat -> Nat) (x : Nat) :
    exists delta : Nat, f (x + delta) = f x := by
  exact Exists.intro 0 rfl

-- versi kuat salah untuk f = identitas, jadi tidak dapat dibuktikan
theorem ada_penyangkal :
    exists x : Nat, exists delta : Nat,
      Not (identitas (x + delta) = identitas x) := by
  exact Exists.intro 0 (Exists.intro 1 (by decide))
\end{verbatim}

Kedua teorema itu lolos tanpa aksioma tambahan, dan justru itulah masalahnya. Teorema pertama
tidak mengatakan apa pun tentang $f$; ia hanya mengatakan bahwa $\delta=0$ adalah gangguan yang
tidak mengubah apa pun, dan itu berlaku untuk setiap fungsi, termasuk pengklasifikasi yang salah
seratus persen. Sertifikat semacam ini lolos semua pemeriksaan dan tidak melindungi apa pun.
Teorema kedua menunjukkan sisi lain: pernyataan yang benar-benar kita inginkan, yaitu versi
berkuantor universal, salah untuk fungsi identitas, sehingga tidak ada bukti yang mungkin
ditemukan.

\begin{tantangan}
Sebuah sertifikat memuat satu baris yang tampak tidak mungkin salah: untuk setiap bilangan asli $a$
dan $b$ berlaku $(a-b)+b=a$. Di atas kertas alasannya sederhana: kurangi, tambahkan lagi, kembali ke
semula. Tulis baris itu di Lean dan tutup dengan \texttt{omega}; pemeriksa menolaknya, dan penolakan
itu bukan karena taktiknya kurang pintar.

Dua pertanyaan. Mengapa pernyataan itu tidak benar untuk semua $a$ dan $b$? Dan syarat tambahan apa
yang membuatnya benar sekaligus tetap dapat diperiksa mesin?

Petunjuk 1: jangan percaya pada tanda kurung; tanyakan apa yang terjadi bila $b$ lebih besar dari
$a$. Petunjuk 2: pengurangan pada bilangan asli berhenti di nol, jadi $0-1=0$, sehingga ruas kiri
menjadi $0-1+1=1$ sedangkan ruas kanan $0$.

Kalau mentok, kembali ke bagian ``Taktik: mengubah tujuan menjadi term'', tempat $n+1$ dan $1+n$
berhenti di bentuk yang berbeda.
\end{tantangan}

Sebelum membaca lebih jauh, coba jawab sendiri. Keluaran mesin uji ini, apa adanya, adalah pembanding
jawabanmu: berkas kedua memuat pernyataan yang sama dengan hipotesis $b\le a$.

\begin{verbatim}
$ lean t.lean
t.lean:1:49: error: omega could not prove the goal:
a possible counterexample may satisfy the constraints
  d >= 0
  c >= 0
  c - d >= 1
where
 c := b
 d := a
$ lean t2.lean
'potong_benar' depends on axioms: [propext, Quot.sound]
\end{verbatim}

Keluaran pertama menyebut kendala $b-a\ge 1$, yaitu penyangkal dengan $b>a$; pada $a=0$ dan $b=1$
ruas kiri bernilai $1$, ruas kanan $0$. Hipotesis $b\le a$ membuat baris yang sama lolos, dengan dua
aksioma standar yang sudah kita kenal dari \texttt{omega}. Bentuk perbaikannya bukan taktik
tambahan, melainkan syarat.

\kertas{Klaim yang boleh dibuat dari bab ini: potongan kode Lean yang ditampilkan telah
dijalankan dengan Lean 4.10.0 tanpa Mathlib pada mesin penulisan, dan keluarannya disalin apa
adanya; keluaran Unicode dicetak ulang dalam ASCII dan pemotongan ditandai dengan tiga titik.
Yang tidak dijalankan: taktik Mathlib (\texttt{linarith}, \texttt{ring}, \texttt{norm\_num},
\texttt{induction'}) dan pustaka Mathlib sendiri, karena tidak terpasang. Yang tidak dibuktikan di
sini dan tidak diklaim: sertifikat apa pun tentang jaringan saraf, model bahasa, atau ketahanan
terhadap serangan. Bab 28 tidak memuat satu pun bukti semacam itu.}

Batas ini ada hubungannya langsung dengan Bab 27 dan Bab 29. Evaluasi model yang jujur menuntut
kita menyebut distribusi tempat angka diukur; verifikasi formal menuntut kita menyebut pernyataan
yang dibuktikan. Keduanya adalah penyakit yang sama dalam bentuk berbeda: klaim yang tidak
menyebut syaratnya. Sertifikat ``model ini terbukti aman'' tanpa pernyataan formal di belakangnya
sama tidak bergunanya dengan ``akurasi 99 persen'' tanpa keterangan data ujinya. Pada Bab 29 dan
Bab 30, ketika kita membahas serangan dan pertahanan, pertanyaan pertama yang harus diajukan
kepada setiap sertifikat ketahanan adalah: berapa kuantornya, dan dalam ruang apa kuantor itu
berjalan?

\section{Contoh hitung}

\tingkat{menengah}

Angka-angka di bagian ini dihitung, bukan diperkirakan. Kolom pertama adalah $n$, kolom kedua
adalah bilangan ganjil ke-$n$ yang ditambahkan pada langkah itu, kolom ketiga adalah $S(n)$, dan
kolom keempat adalah $n^{2}$.

\begin{center}
\begin{tabular}{rrrr}
\hline
$n$ & suku ke-$n$ & $S(n)$ & $n^{2}$ \\
\hline
0  & --  & 0     & 0 \\
1  & 1   & 1     & 1 \\
2  & 3   & 4     & 4 \\
5  & 9   & 25    & 25 \\
10 & 19  & 100   & 100 \\
20 & 39  & 400   & 400 \\
100 & 199 & 10000 & 10000 \\
\hline
\end{tabular}
\end{center}

Nilai $S(100)=10000$ patut diperiksa dengan tangan sekali, karena ia menunjukkan bahwa rumus itu
bukan kebetulan untuk bilangan kecil: jumlah seratus bilangan ganjil, dari $1$ sampai $199$,
sama dengan $100\cdot 100$. Pemeriksaan dengan Python atas seluruh nilai $n<5000$ memberi
\texttt{True}, dan pemeriksaan identitas aljabar $(k+1)^{2}=k^{2}+2k+1$ untuk seluruh $k<100000$
juga memberi \texttt{True}. Kedua pemeriksaan itu bukan bukti, tetapi keduanya menutup kemungkinan
bahwa rumus yang dibuktikan salah ketik.

Ada satu angka lagi yang lebih informatif daripada tabel di atas, yaitu ukuran bukti setelah
dielaborasi. Kalimat taktik untuk \texttt{oddSum\_sq} panjangnya delapan baris. Ketika Lean
mencetak term hasil elaborasinya dengan \texttt{\#print}, hasilnya 10.783 karakter dalam 149
baris. Artinya, satu baris taktik rata-rata mengembang menjadi sekitar sembilan belas baris term.
Tidak ada manusia yang membaca term itu; yang membacanya adalah inti, dan itulah pembagian kerja
yang menjadi inti seluruh metode ini.

Angka terakhir menyangkut batas pemeriksaan langsung. Pada pengaturan bawaan, batas kedalaman
rekursi bernilai 1000. Dengan batas itu, taktik \texttt{decide} berhasil membuktikan
$S(n)=n^{2}$ hanya untuk $n\le 122$. Percobaan pertama yang gagal adalah $n=123$, dan pesan yang
keluar menyebut bahwa kedalaman rekursi maksimum telah tercapai. Bila batas dinaikkan ke 100000,
nilai $n=2000$ berhasil dibuktikan dalam 0,48 detik. Pada $n=20000$ pemeriksaan gagal karena
rekursi terlalu dalam untuk ruang tumpukan. Jadi pemeriksaan langsung bukan alternatif yang
berskala, sedangkan bukti induksi tidak peduli pada besarnya $n$.

\section{Lab komputasi}

\tingkat{lanjut}

Lab ini menjalankan dua hal: berkas Lean yang sama dengan yang dibahas di atas, dan pemeriksaan
silang dengan Python. Berkas Lean ditulis tanpa \texttt{import}, sehingga hanya memakai inti.

\begin{lab}
\begin{verbatim}
-- BilanganGanjil.lean : Lean 4.10.0, tanpa Mathlib
def oddSum : Nat -> Nat
  | 0     => 0
  | n + 1 => oddSum n + (2 * n + 1)

theorem oddSum_sq (n : Nat) : oddSum n = n * n := by
  induction n with
  | zero => rfl
  | succ k ih =>
    rw [oddSum]
    rw [ih]
    simp [Nat.succ_eq_add_one, Nat.mul_succ, Nat.add_mul, Nat.one_mul]
    omega

#eval oddSum 0
#eval oddSum 5
#eval oddSum 10
#eval oddSum 100
#print axioms oddSum_sq
\end{verbatim}

\begin{verbatim}
def odd_sum(n):
    return sum(2 * i + 1 for i in range(n))

print([odd_sum(n) for n in range(11)])
print(odd_sum(20), odd_sum(100))
print(all(odd_sum(n) == n * n for n in range(5000)))
print(all((k + 1) * (k + 1) == k * k + 2 * k + 1 for k in range(100000)))
\end{verbatim}
\end{lab}

Perintah dan keluaran, apa adanya:

\begin{verbatim}
$ lean --version
Lean (version 4.10.0, x86_64-unknown-linux-gnu, commit c375e19f6b65, Release)
$ time lean BilanganGanjil.lean
0
25
100
10000
'oddSum_sq' depends on axioms: [propext, Quot.sound]
waktu 0.14 s
$ python3 BilanganGanjil.py
[0, 1, 4, 9, 16, 25, 36, 49, 64, 81, 100]
400 10000
True
True
\end{verbatim}

Empat baris pertama keluaran Lean adalah hasil \texttt{\#eval}, dan keempatnya cocok dengan tabel
di bagian sebelumnya. Baris kelima adalah daftar aksioma, dan isinya perlu dibaca dengan hati-hati:
\texttt{propext} dan \texttt{Quot.sound} bukan kesalahan, melainkan akibat dari cara
\texttt{omega} bekerja. Keluaran Python menguatkan hal yang sama dari arah lain: daftar sebelas
nilai pertama adalah kuadrat sempurna, $S(100)=10000$, dan kedua pemeriksaan menyeluruh mencetak
\texttt{True}. Waktu 0,14 detik adalah waktu mesin uji ini, dan angka semacam itu bergantung pada
perangkat keras; yang bermakna bukan angkanya, melainkan faktanya bahwa pemeriksaan selesai
sebelum kita sempat membaca berkasnya.

Pembaca yang ingin mengulang lab ini memerlukan dua hal. Pertama, Lean 4 yang dapat dipasang
lewat \texttt{elan}, yaitu pengelola versi resmi proyek Lean; setelah pemasangan, perintah
\texttt{lean --version} harus mencetak versi. Kedua, berkas di atas disimpan sebagai
\texttt{BilanganGanjil.lean} dan dijalankan dengan \texttt{lean BilanganGanjil.lean} dari direktori
yang sama; tidak perlu proyek \texttt{lake} karena tidak ada impor. Untuk memakai \texttt{linarith}
atau \texttt{ring}, pembaca perlu proyek dengan Mathlib, dan sejak saat itu \texttt{lean} saja
tidak cukup; perintahnya menjadi \texttt{lake env lean}. Perbedaan ini sering menjadi sumber
kebingungan pertama bagi pemula, dan menyebutkannya lebih baik daripada membiarkannya ditemukan
lewat pesan \texttt{unknown tactic}.

\begin{resep}
\begin{enumerate}
  \item Buat berkas teks biasa bernama \texttt{Uji.lean}; ekstensi itulah yang membuat perintah
        \texttt{lean} mengenalinya sebagai berkas sumber.
  \item Tulis semua \texttt{import} di baris paling atas, satu per baris, sebelum deklarasi apa pun.
  \item Tulis pernyataan lebih dulu, lengkap dengan tipe setiap variabelnya, lalu bukti sesudah
        \texttt{:= by}.
  \item Jalankan \texttt{lean Uji.lean} dari direktori berkas itu dan baca keluaran dari baris
        pertama; kesalahan pertama biasanya yang sebenarnya.
  \item Tanyakan \texttt{\#print axioms nama\_teorema} sebelum menyebut hasilnya terbukti.
\end{enumerate}
Peringatan urutan: \texttt{import} tidak boleh didahului deklarasi. Satu \texttt{def} yang
diselipkan di atasnya membuat Lean menolak seluruh berkas dengan pesan
\texttt{invalid 'import' command}, dan berkas itu tidak diperiksa sama sekali. Urutan yang benar
bukan selera penulis; ia syarat agar berkas dapat diperiksa.
\end{resep}

\begin{untukdosen}
Urutan mengajar yang berhasil di kelas: mulai dari kegagalan \texttt{rfl} pada $n+1=1+n$, bukan
dari definisi. Mahasiswa yang melihat bahwa hal yang ``jelas'' ditolak akan bertanya, dan
pertanyaan itu adalah pintu masuk ke Curry--Howard. Setelah itu baru definisi proposisi sebagai
tipe, lalu tiga taktik pertama, lalu induksi. Sisakan satu sesi penuh untuk \texttt{\#print
axioms}; di situlah mahasiswa berubah dari pengguna menjadi pemeriksa.

Dua miskonsepsi yang paling sering muncul. Pertama, ``Lean membuktikan bahwa program saya benar''.
Yang benar: Lean memeriksa bahwa term bertipe pernyataan tertentu. Bila pernyataannya tidak
menyebut hal yang penting, misalnya kuantor atas gangguan, sertifikatnya kosong. Kedua,
``\texttt{simp} adalah bukti''. \texttt{simp} adalah prosedur yang keluarannya bergantung pada
kumpulan lemma aktif, sehingga bukti yang bergantung padanya dapat berhenti bekerja ketika pustaka
berubah; karena itu latih mahasiswa menulis daftar lemma secara eksplisit.

Aktivitas papan tulis yang berhasil: minta dua kelompok menulis spesifikasi untuk ``pengklasifikasi
yang tahan gangguan''. Satu kelompok hampir selalu menulis kuantor eksistensial. Tempelkan kedua
rumusan di papan, lalu buktikan yang lemah untuk fungsi konstan. Kelas akan melihat sendiri bahwa
kalimat yang terbukti dan kalimat yang dimaksud adalah dua kalimat berbeda.

Tanda kesiapan lanjut: mahasiswa dapat menjelaskan mengapa $n+1=1+n$ memerlukan lemma, dapat
menulis bukti induksi tanpa menyalin, dan dapat menyebut satu aksioma yang dipakai teoremanya
sendiri.
\end{untukdosen}

\section{Ringkasan}

\begin{ringkasan}
\begin{itemize}
  \item Bukti informal meyakinkan pembaca terlatih; bukti formal adalah term yang tipenya adalah
        pernyataannya, diperiksa inti yang kecil dan tidak tahu matematika.
  \item Menurut Curry--Howard, implikasi adalah tipe fungsi, konjungsi adalah pasangan, kuantor
        universal adalah fungsi dependen, dan $\bot$ tidak berpenghuni.
  \item Taktik mengubah keadaan tujuan; \texttt{intro}, \texttt{exact}, \texttt{apply},
        \texttt{rw}, \texttt{cases}, dan \texttt{induction} cukup untuk aritmetika dasar.
  \item Induksi membuktikan pernyataan untuk semua $n$ dengan dua tujuan, bukan dengan memeriksa
        semua $n$; \texttt{decide} sebaliknya berhenti berskala sangat cepat.
  \item \texttt{\#print axioms} adalah alat pemeriksaan utama. Taktik \texttt{native\_decide}
        meninggalkan aksioma \texttt{Lean.ofReduceBool}. Taktik \texttt{sorry} meninggalkan
        \texttt{sorryAx}. Aksioma buatan sendiri meninggalkan namanya sendiri.
  \item Verifikasi menjamin term cocok dengan pernyataan, bukan bahwa pernyataan cocok dengan
        maksud; spesifikasi lemah menghasilkan sertifikat yang kosong.
\end{itemize}
\end{ringkasan}

\section{Latihan}

\begin{latihan}
\begin{enumerate}
  \item Nilai klaim ini: \emph{bila sebuah berkas Lean diterima tanpa pesan kesalahan, maka semua
        pernyataan di dalamnya benar}. Tentukan benar, salah, atau belum cukup bukti, dan sebutkan
        satu perintah yang dapat dipakai untuk memeriksa klaim itu pada berkas tertentu.

  \item Buktikan di Lean, tanpa Mathlib, bahwa untuk setiap $n$ berlaku $2\cdot S(n)=n\cdot(2n)$
        dengan $S$ seperti di bab ini. Nyatakan lebih dulu lemma bantu yang kamu perlukan.

  \item Buktikan bahwa jumlah $n$ bilangan genap pertama, yaitu $0+2+\dots+(2n-2)$, sama dengan
        $n(n-1)$. Tulis bukti informalnya, lalu tulis versi Lean-nya dan sebutkan taktik yang
        dipakai pada setiap langkah.

  \item Hitung dengan Python: jumlah 37 bilangan ganjil pertama, jumlah 250 bilangan ganjil
        pertama, dan periksa bahwa keduanya kuadrat sempurna. Sertakan kode dan keluarannya.

  \item Hitung dengan Python banyaknya kasus yang harus diperiksa bila seseorang ingin membuktikan
        $S(n)=n^{2}$ dengan pemeriksaan langsung untuk semua $n\le 1000$, lalu bandingkan dengan
        banyaknya tujuan yang dibangkitkan bukti induksi.

  \item Jelaskan mengapa \texttt{rfl} menutup tujuan $n+0=n$ tetapi tidak menutup $0+n=n$. Petunjuk
        ada pada cara penjumlahan bilangan asli didefinisikan secara rekursif; sebutkan argumen
        mana yang disusutkan lebih dulu.

  \item Tulis spesifikasi lemah dan spesifikasi kuat untuk klaim ``pengklasifikasi ini tidak
        berubah pada gangguan berukuran paling banyak satu''. Buktikan versi lemahnya di Lean
        untuk fungsi apa pun, lalu jelaskan mengapa versi kuatnya tidak dapat dibuktikan untuk
        fungsi identitas.

  \item Periksa dengan \texttt{\#print axioms} tiga teorema yang kamu tulis sendiri: satu memakai
        \texttt{decide}, satu memakai \texttt{omega}, satu memakai \texttt{sorry}. Laporkan
        ketiga daftar aksiomanya dan terangkan asal setiap nama.

  \item Sebuah sertifikat memuat \texttt{theorem aman : exists delta : Nat, f (x + delta) = f x}
        dengan \texttt{delta} sebagai variabel bebas di dalam konteks. Jelaskan dua kesalahan
        yang mungkin ada pada pernyataan itu, dan tulis versi yang kamu anggap benar.

  \item Buktikan di Lean bahwa $\neg(P\wedge\neg P)$ untuk setiap proposisi $P$. Petunjuk: mulai
        dengan \texttt{intro h}, lalu ambil komponennya dengan \texttt{h.1} dan \texttt{h.2}, dan
        gunakan komponen kedua pada komponen pertama.
\end{enumerate}

Petunjuk untuk tiga butir tersulit. Butir 2: lemma bantu yang diperlukan berbentuk
$2\,(k\cdot k+(2k+1))=(k+1)\cdot(2(k+1))$, dan setelah \texttt{rw} pada definisi serta hipotesis
induksi, sisa tujuannya linear sehingga \texttt{omega} cukup. Butir 7: pada versi lemah, saksi
$\delta=0$ membuat kedua sisi sama secara definisi, jadi \texttt{rfl} langsung menutup; pada versi
kuat, fungsi identitas memberi penyangkal dengan $\delta=1$. Butir 9: perhatikan bahwa
\texttt{delta} yang bebas dapat dipilih oleh pembaca pernyataan, bukan oleh penyerang, sehingga
pernyataan itu tidak menyatakan ketahanan apa pun; kesalahan kedua adalah tidak adanya kuantor
atas $x$.
\end{latihan}

\begin{untukpeneliti}
Yang benar-benar terbukti di bab ini hanyalah pernyataan aritmetika kecil: $S(n)=n^{2}$ untuk
semua bilangan asli, beberapa tautologi proposisional, dan dua pernyataan tentang spesifikasi yang
terlalu lemah. Semuanya diperiksa inti Lean 4.10.0 tanpa Mathlib, dan daftar aksiomanya
dilaporkan. Yang dipopulerkan jauh melampaui itu: frasa ``terverifikasi secara formal'' sering
dipakai untuk ``diperiksa alat otomatis'' atau bahkan ``diuji menyeluruh''. Hanya yang pertama
menghasilkan objek yang dapat diperiksa ulang pihak ketiga.

Sumber primer. Avigad, de Moura, Kong, dan Ullrich, \emph{Theorem Proving in Lean 4}, bagian
``Propositions and Proofs'' (proposisi sebagai tipe, irrelevansi bukti, peringatan aksioma),
``Tactics'' (keadaan tujuan), dan ``Induction and Recursion'' (\texttt{induction},
\texttt{sorryAx}). Avigad dan Massot, \emph{Mathematics in Lean}, Bab 2, 3.1, 5.2, dan 8; edisi
bahasa Indonesia memakai penomoran bagian yang sama.

Pertanyaan terbuka. Bab ini tidak memuat bukti formal tentang jaringan saraf. Spesifikasi ketahanan yang menarik
berbentuk
$\forall x\,\forall\delta,\ \lVert\delta\rVert\le\varepsilon\to f(x+\delta)=f(x)$, yaitu
pernyataan atas ruang kontinu dengan fungsi berparameter ratusan ribu. Memformalkannya menuntut
aritmetika real yang dapat diputuskan pada skala itu, dan pustaka Lean sekarang tidak menyediakannya
untuk model terlatih. Klaim ``kami memverifikasi ketahanan model'' di makalah
biasanya merujuk verifikasi dengan abstraksi interval atau SMT, bukan bukti Lean; keduanya sah,
tetapi jaminannya berbeda dan tidak boleh disebut dengan kata yang sama.

Batas keberlakuan. Semua klaim di bab ini berlaku untuk Lean 4.10.0 dengan inti bawaannya, dan
rincian keluarannya dapat berubah antar versi. Yang tidak berubah adalah struktur argumennya:
pernyataan, term, inti, daftar aksioma.
\end{untukpeneliti}

\peta{Bab ini mengikuti Avigad, de Moura, Kong, dan Ullrich, \emph{Theorem Proving in Lean 4},
bagian ``Propositions and Proofs'' (proposisi sebagai tipe, Curry--Howard, irrelevansi bukti,
peringatan aksioma), ``Quantifiers and Equality'', ``Tactics'' (keadaan tujuan, \texttt{intro},
\texttt{apply}, \texttt{exact}), dan ``Induction and Recursion'' (\texttt{induction},
\texttt{sorry} dan \texttt{sorryAx}). Padanan dalam Avigad dan Massot, \emph{Mathematics in Lean},
ada pada Bab 2 (Dasar-dasar, khususnya 2.3 Menggunakan Teorema dan Lema serta 2.4 Lebih Banyak
Contoh Menggunakan apply dan rw), Bab 3.1 (Implikasi dan Kuantor Universal), Bab 5.2 (Induksi dan
Rekursi), dan Bab 8 (Hierarki) untuk typeclass; edisi bahasa Indonesia mempertahankan penomoran
bagian yang sama. Untuk kaitan dengan praktik evaluasi, lihat Bab 27 buku ini; untuk serangan yang
membuat spesifikasi ketahanan menjadi tidak trivial, lihat Bab 27 dan Bab 28.}
